\documentclass[11pt]{article}
\usepackage[a4paper,margin=1in]{geometry}
\usepackage[T1]{fontenc}
\usepackage[utf8]{inputenc}
\usepackage{lmodern}
\usepackage{amsmath,amssymb}
\usepackage{booktabs}
\usepackage{hyperref}

\title{CP-SAT Tearing: Formulation, Execution, and Auditability}
\author{Project +Digital}
\date{\today}

\begin{document}
\maketitle

\section{Purpose}
This repository provides a reproducible CP-SAT pipeline for structural variable
classification in equation systems with closed tearing semantics. The goals are:
\begin{itemize}
  \item run all studied cases under a deterministic configuration;
  \item classify variables into operationally meaningful structural classes;
  \item generate auditable artifacts for traceability and reporting.
\end{itemize}

\section{Modeling summary}
Given equations $Q$ and variables $V$, each equation selects at most one output
variable. Local tear cuts may be introduced to break cycles in the execution
graph. The model optimizes structural coverage in phases, with deterministic
settings and explicit status reporting.

\subsection{Main sets}
\begin{itemize}
  \item Measured variables ($V^{m}$): directly instrumented variables.
  \item Known constants ($V^{k}$): known parameters available as inputs.
  \item Directly inferred variables: acyclic propagation from known inputs.
  \item Closed-loop inferred variables: inferred in cyclic blocks resolved by
  closed tearing.
  \item Open tears and open-conditioned variables: variables that require
  external tear values to become numerically available.
\end{itemize}

\subsection{Coverage metrics}
The implementation reports three complementary metrics:
\begin{align}
C_{dir} &= \frac{|V^{m} \cup V^{k} \cup V^{dir}|}{|V|}, \\
C_{cl}  &= \frac{|V^{m} \cup V^{k} \cup V^{dir} \cup V^{loop} \cup V^{tear,closed}|}{|V|}, \\
C_{ext} &= \frac{|V^{m} \cup V^{k} \cup V^{dir} \cup V^{loop} \cup V^{tear,closed} \cup V^{cond,open} \cup V^{tear,closed,ext}|}{|V|}.
\end{align}

Interpretation:
\begin{itemize}
  \item $C_{dir}$: direct acyclic reach;
  \item $C_{cl}$: autonomous closed coverage (no open-tear dependency);
  \item $C_{ext}$: conditional reach if open tears are supplied externally.
\end{itemize}

\section{Repository layout}
\begin{itemize}
  \item \texttt{src/}: solver core, case registry, and case definitions.
  \item \texttt{scripts/}: execution entrypoints.
  \item \texttt{tests/}: semantic smoke tests.
  \item \texttt{results/}: generated artifacts (ignored by git).
\end{itemize}

\section{How to run}
\subsection{Environment setup}
\begin{verbatim}
python3 -m pip install -r requirements.txt
\end{verbatim}

\subsection{Semantic checks}
\begin{verbatim}
python3 tests/test_tearing_semantics.py
\end{verbatim}

\subsection{Run studied cases}
\begin{verbatim}
python3 scripts/run_cases.py
\end{verbatim}

Optional controls:
\begin{verbatim}
python3 scripts/run_cases.py --time-limit 60 --allow-feasible
python3 scripts/run_cases.py --case 01_urs_ideal --case 18_narasimhan
\end{verbatim}

\subsection{Additional analyses}
\begin{verbatim}
python3 scripts/run_matching_ablation.py
python3 scripts/run_time_benchmark.py --case 18_narasimhan --k 17 --n-samples 1000
\end{verbatim}

\section{Audit trail}
Each run creates a timestamped folder:
\begin{verbatim}
results/run_YYYYMMDD_HHMMSS/
\end{verbatim}
with consolidated artifacts:
\begin{itemize}
  \item \texttt{tearing\_results\_table.csv}
  \item \texttt{tearing\_results\_table.xlsx}
  \item \texttt{tearing\_results\_table.tex}
  \item \texttt{auditoria.json}
\end{itemize}

The \texttt{auditoria.json} file records:
\begin{itemize}
  \item semantics identifier;
  \item execution timestamp, Python version, platform;
  \item solver configuration;
  \item global checks (all optimal, DAG consistency, etc.).
\end{itemize}

\section{Determinism notes}
The solver is configured with fixed seed and deterministic phase flow. Minor
runtime variation is expected across machines, but structural classifications
and coverage metrics should remain stable under equivalent dependencies.

\end{document}
